\documentclass[12pt]{letter}
\usepackage{geometry}
\geometry{margin=1in}
\usepackage{hyperref}

\begin{document}

\begin{letter}{%
Editor-in-Chief\\
\emph{Artificial Intelligence and Law}\\
Springer Nature}

\opening{Dear Editor,}

I am pleased to submit the manuscript ``Algorithmic Jurisprudence: The Topology of Law, Constitutionality as Homology, and Legal Reasoning as Optimal Pathfinding'' for consideration as an Original Research article in \emph{Artificial Intelligence and Law}.

The paper proposes a mathematical framework in which the space of legal states is a directed weighted simplicial complex (the \emph{judicial complex}), legal disputes are optimal pathfinding problems, and constitutionality is a topological invariant computable via path homology on directed graphs.  The framework extends the Geometric Ethics programme from ethics to law, introducing two legal invariance principles (LIP, JBIP), proving that equal protection is a gauge symmetry, that due process is quotient regularity, that liability--damages are conserved in closed bilateral disputes, and that constitutionality is equivalent to preservation of the path homology of the constitutional subcomplex.

The paper makes seven specific contributions:
\begin{enumerate}
    \item Construction of the judicial complex with Mahalanobis-weighted directed edges and regime filtration;
    \item Legal invariance principles extending the Bond Invariance Principle to Hohfeldian legal structure;
    \item The octahedral gauge group $D_4 \rtimes D_4$ for the full Hohfeldian octad, with Wilson loop consistency tests;
    \item A topological constitutionality criterion via path homology (Grigor'yan--Lin--Muranov--Yau theory);
    \item Liability--damages conservation derived from Hohfeldian correlative structure;
    \item Formalization of litigation as $A^*$ search with doctrinal heuristics;
    \item A concrete NLP calibration pipeline (LaBSE embedding $\to$ linear probes $\to$ Mahalanobis edge weights) with validation safeguards.
\end{enumerate}

This work is relevant to the journal's scope in three ways.  First, it provides formal computational models of legal knowledge and reasoning---the judicial complex is a constructible data structure from case databases and citation networks.  Second, it directly addresses the implications of AI in law by proposing the Legal Bond Index as a quantitative test for algorithmic fairness in legal AI systems.  Third, it bridges two literatures central to the journal: Hohfeldian formal analysis and computational legal reasoning, unifying them through gauge theory.

The manuscript has not been published previously and is not under consideration elsewhere.  All data and code from the parent framework are publicly available.

I have suggested six reviewers with expertise spanning computational law, Hohfeldian analysis, legal argumentation, and AI and law, from six different countries and institutions.

Thank you for considering this submission.

\closing{Sincerely,\\[2em]
Andrew H. Bond\\
Senior Lecturer, Department of Computer Engineering\\
San Jos\'{e} State University\\
\href{mailto:andrew.bond@sjsu.edu}{andrew.bond@sjsu.edu}\\
ORCID: 0009-0003-2599-6158}

\end{letter}
\end{document}
